% page 389 du fac-simile (image p-388 du scan)
% bloc 170.2 x 242.0 mm ; marge gauche 31.8 mm
\pgc{10.160mm}{0.000mm}{
\l{\cel{57}381}
\l{}
\l{\cel{4}monodroma. (\sou{matem.)}\cel{2}Funciono monodroma, interne di areo do-}
\l{\cel{7}nita. Funciono\cel{1}\sou{F(z)}\cel{1}esas monodroma interne di areo S, kande,}
\l{\cel{7}la punto z variante interne di la areo, la funciono\cel{1}\sou{F(z)}}
\l{\cel{7}retroprenas sempre la sama valoro en sama punto. - EFIS.}
\l{}
\l{\cel{4}\sou{monogama}. I. Qua esas mariajita kun nur un spozo. - II. Qua}
\l{\cel{7}spozeskis nur unfoye. - DEFIRS.}
\l{}
\l{\cel{4}\sou{monogena}. (\sou{biol.)}\cel{3}Qua agas ta sorto di genito en qua la enti}
\l{\cel{7}vivanta su reproduktas nemediate, e kun fazi inter-identa di}
\l{\cel{7}su-developo, per ovo o per ovuli, opoze a la digeneso (ge-}
\l{\cel{7}nita alternanta). -\cel{2}DEFIS.}
\l{}
\l{\cel{4}\sou{monografio}. Studiuro pri punto specala di historio, cienco,}
\l{\cel{7}e c. - DEFIRS.}
\l{}
\l{\cel{4}\sou{monogramo}. (\sou{skribarto}).I. Kombinuro ek plura literi, unionita}
\l{\cel{7}tale ke un streko, un slingo, uzesas por du o tri literi}
\l{\cel{7}interdiferanta.\cel{2}- II. Kombinuro de literi quin ula artisti}
\l{\cel{7}stampas sur sua verki - ula kolekteri, sur la peci di lia}
\l{\cel{7}kolektaji. - DEFIRS.}
\l{}
\l{\cel{4}\sou{monoika}. (\sou{bot.)}\cel{2}Qua portas, sur un stipo, floruli e florini. -}
\l{}
\l{}
\l{\cel{4}\sou{monoklo}. Mikra okulvitruno, por un ek la okuli. - DEFIRS.}
\l{}
\l{\cel{4}\sou{monokroma.}\cel{2}Qua esas un-kolora. - DEFIS.}
\l{}
\l{\cel{4}\sou{monoklina}.\cel{2}(\sou{bot.)}\cel{2}En qua la du sexui un-esas en la sama}
\l{\cel{7}floro. -}
\l{}
\l{\cel{4}\sou{monolito}. Monumento (kolono, obelisko, tombo, e c.) ek nur un}
\l{\cel{7}petro. - DEFIRS.}
\l{}
\l{\cel{4}\sou{monomio}. (\sou{algebro}). Expresuro inter la parti di qua ne existas}
\l{\cel{7}signo di adiciono o di sustraciono. DEFIS.}
\l{}
\l{\cel{4}\sou{monologo}. I. Ceno teatrala en qua un persono - qua esas, o}
\l{\cel{7}qua kredas esar, sola - parolas a su ipsa. - II. Diskurso}
\l{\cel{7}da persono qua ne lasas la kun-esanti parolar liafoye.-}
\l{\cel{7}DEFIRS.}
\l{}
\l{\cel{4}monopolo. I. Privilejo exkluziva pri vendar ulo. - II. Yuro}
\l{\cel{7}exkluziva, posedata da mikra quanto de homi. - DEFIRS.}
\l{}
\l{\cel{4}monosperma. (\sou{bot.)}\cel{3}Di qua la frukto kontenas nur un grano}
\l{\cel{7}(semino). - DEFIS.}
\l{}
\l{\cel{4}monoteismo. Kredo ye deo unika. - DEFIRS.}
\l{}
\l{\cel{4}monotona. I. Qua esas sempre sama-tona. Qua fatigas pro uzar}
\l{\cel{7}konstante la sama tono. - II.Qua fatigas pro la repeto di}
\l{\cel{7}la sama kozi, pro la itero di la sama ago. - DEFIRS.}
}
